% Generated draft: requires mathematical and source review.
\documentclass[11pt]{article}
\usepackage[margin=25mm]{geometry}
\usepackage{fontspec}
\setmainfont{Latin Modern Roman}
\usepackage{amsmath,amssymb,mathtools,unicode-math}
\setmathfont{Latin Modern Math}
\usepackage{enumitem,needspace,xurl,hyperref}
\setlength{\parindent}{0pt}
\setlength{\parskip}{5pt}
\setlength{\emergencystretch}{3em}
\newcommand{\N}{\mathbb{N}}
\newcommand{\Z}{\mathbb{Z}}
\newcommand{\Q}{\mathbb{Q}}
\newcommand{\R}{\mathbb{R}}
\newcommand{\C}{\mathbb{C}}
\newcommand{\F}{\mathbb{F}}
\newcommand{\A}{\mathbb{A}}
\newcommand{\PP}{\mathbb{P}}
\newcommand{\E}{\mathbb{E}}
\newcommand{\eps}{\varepsilon}
\newcommand{\dd}{\,\mathrm{d}}
\begin{document}
\section*{UnsolvedMath ID: 30004571}
\textbf{Source problem code:} catalogue entry
\subsection{Definitions and mathematical statement}
Let \(\Omega=(0,1)\). The supplied problem concerns a spatial profile \(f\), a time-dependent source factor \(h\), and a temperature \(u\) satisfying
\[
\begin{cases}
\partial_t u(t,x)-\partial_x^2u(t,x)=h(t)f(x), & x\in(0,1),\ t>0,\\
u(t,0)=u(t,1)=0, & t>0,\\
u(0,x)=0, & x\in(0,1).
\end{cases}
\]
The measurement at a fixed observation point is \(m(t)=u(t,x_0)\).

A conditional formulation of the intended uniqueness question is as follows. Fix an observation horizon \(T>0\), a prescribed profile \(f\), an admissible class \(\mathcal H\) of time-dependent factors, and a solution framework in which each \(h\in\mathcal H\) has a unique solution \(u_h\) and the point trace exists. For \(x_0\in(0,1)\), define
\[
\mathcal M_{f,x_0,T}(h)=u_h(\cdot,x_0)\big|_{(0,T)}.
\]
Determine conditions on \(f\), \(\mathcal H\), the solution framework, and \(x_0\) under which this map is injective; in particular, characterize the observation locations for which
\[
\forall h_1,h_2\in\mathcal H,\qquad
\mathcal M_{f,x_0,T}(h_1)=\mathcal M_{f,x_0,T}(h_2)
\ \Longrightarrow\ h_1=h_2\quad\text{on }(0,T).
\]
Equality here must be interpreted in the chosen function spaces.

This is a formulation template requiring additional hypotheses, rather than a fully specified theorem or conjecture. The record does not specify the observation duration, regularity or admissible classes, solution concept, or nondegeneracy assumptions on \(f\), and does not explicitly state that \(f\) is known. Treating \(f\) as prescribed and \(x_0\) as an interior point expresses the intended recovery question but requires confirmation. If boundary observation points are allowed, their traces are already fixed by the homogeneous Dirichlet conditions.

The supplied dated assessment describes conditional source recovery and nonuniqueness associated with eigenfunction zeros, but gives no precise hypotheses or quantifiers for either assertion. These qualifications are retained as supplied context, not asserted here as verified results.
\subsection{Short English statement}
Can the time factor \(h\) in a separated heat source \(h(t)f(x)\) be uniquely recovered from the temperature history at one spatial point, and which observation points permit recovery? A precise answer depends on the spatial profile, observation interval, admissible functions, and solution framework, which the record leaves unspecified. The supplied status is \texttt{partially\_solved}; this draft does not independently verify that assessment.
\subsection{Sources}
\begin{itemize}
\item Supplied local UnsolvedMath registry record, permanent ID \(30004571\), title: \emph{Single-Point Recovery of Time-Dependent Heat Sources}; exact source problem code: \texttt{catalogue entry}. Project registry paths are \texttt{data/unsolvedmath/problems.json} and the supplementary \texttt{catalog/problems.json}; the supplied extract does not identify which path contains this record.
\item Supplied report citation: \emph{Computational Inverse Problems for Partial Differential Equations} (2021), Oberwolfach Reports, pp.\ 1903--1954, DOI 10.4171/owr/2020/39. Literature-source label: ``OWR 2020/39 report.'' \url{https://doi.org/10.4171/owr/2020/39}. The record gives proposed year 2020 and names no individual proposer. Its background attributes the collection to Oberwolfach workshop reports and their open-problem discussions.
\item Supplied literature-source label: ``Hasanov and Slodicka, inverse source from point data.'' \url{https://doi.org/10.1016/j.aml.2012.08.013}.
\item Supplied literature-source label: ``Harnack inequality for fractional diffusion.'' \url{https://www.sciencedirect.com/science/article/pii/S002203961730013X}.
\item Literature checked at: \texttt{2026-08-22}. This date records supplied dated curation, not an independently verified assessment. The supplied status is \texttt{partially\_solved}, with the qualification that recovery requires nondegeneracy and observation-location assumptions and that eigenfunction zeros can cause nonuniqueness. The record explicitly calls this a dated literature triage derived from supplied OWR curation and requests independent verification of time-sensitive status and cited claims. No cited source was visited for this draft.
\end{itemize}
\end{document}
